% ============================================================
%  grammar.tex  —  Polish Grammar Foundations
%  Referenced by polish_study.tex via % ============================================================
%  grammar.tex  —  Polish Grammar Foundations
%  Referenced by polish_study.tex via % ============================================================
%  grammar.tex  —  Polish Grammar Foundations
%  Referenced by polish_study.tex via % ============================================================
%  grammar.tex  —  Polish Grammar Foundations
%  Referenced by polish_study.tex via \input{grammar}
% ============================================================

\chapter*{Polish Grammar Foundations}
\addcontentsline{toc}{chapter}{Polish Grammar Foundations}
\markboth{Polish Grammar Foundations}{}

This chapter provides a compact reference for the grammatical structures 
encountered throughout the reading. It is not intended as a complete grammar 
— rather, as a touchstone for constructs that recur in the text. Cross-references 
to passages where each construct first appears are noted where relevant.

% ────────────────────────────────────────────────────────────
\section{Nouns and Gender}

Polish nouns have three grammatical genders: \textbf{masculine}, \textbf{feminine}, 
and \textbf{neuter}. Unlike French, where gender must largely be memorised, Polish 
gender is often inferable from the noun's ending.

\begin{description}
  \item[Masculine] Typically end in a consonant: \polword{grób}{grave}, 
    \polword{pogrzeb}{funeral}, \polword{syn}{son}.
  \item[Feminine] Typically end in \textbf{-a} or \textbf{-i/y}: 
    \polword{wdowa}{widow}, \polword{grupka}{small group}, 
    \polword{nadzieja}{hope}.
  \item[Neuter] Typically end in \textbf{-o} or \textbf{-e}: 
    \polword{dzieciństwo}{childhood}, \polword{dziecko}{child}.
\end{description}

\grambox{Note on animacy}{Masculine nouns are further divided into 
\textit{animate} (persons and animals) and \textit{inanimate}. This distinction 
affects the accusative case ending and is particularly important for direct objects.}

Nouns also form diminutives extensively in Polish, often with suffixes 
\textbf{-ka}, \textbf{-ko}, \textbf{-ek}, \textbf{-eczka}. These indicate 
smallness but frequently carry emotional colouring — affection, sympathy, or 
irony — rather than literal size. Examples encountered in Chapter 1: 
\polword{grupka}{little group} (from \textit{grupa}), 
\polword{garstka}{small handful} (from \textit{garść}), 
\polword{flaszka}{bottle, colloquial} (from \textit{flasza}),
\polword{podstawówka}{primary school, affectionate} (from \textit{szkoła podstawowa}).

% ────────────────────────────────────────────────────────────
\section{The Seven Cases}

Polish has seven grammatical cases. Each case encodes a grammatical 
relationship that English expresses through word order or prepositions. 
The Latin equivalents are noted as a reference point.

\begin{description}

  \item[Nominative (\textit{mianownik})] The subject of the sentence. 
    Who or what performs the action. Equivalent to Latin nominative.\\
    \textit{Groszek \textbf{powiedział}\ldots} — Groszek said\ldots

  \item[Genitive (\textit{dopełniacz})] Possession, absence, quantity, 
    and the object of many prepositions (\textit{do, od, z, bez, dla, oprócz}). 
    The most frequently used case after the nominative. Equivalent to Latin genitive.\\
    \textit{dla \textbf{uczczenia} jego \textbf{pamięci}} — for the honouring of his memory\\
    \textit{garstki \textbf{obecnych}} — a handful of those present

  \item[Dative (\textit{celownik})] The indirect object — to whom or for whom. 
    Equivalent to Latin dative. Encodes dedication, as seen in the book's 
    epigraph: \textit{\textbf{Ojcu}, \textbf{Synowi}} — to [my] father, to [my] son.

  \item[Accusative (\textit{biernik})] The direct object of most verbs, 
    and the object of motion prepositions (\textit{na, w, przez, za}).
    Equivalent to Latin accusative.\\
    \textit{Zostaw \textbf{tę flaszkę}!} — Leave that bottle!

  \item[Instrumental (\textit{narzędnik})] Means, accompaniment, and 
    identity (used after the verb \textit{być} — to be, for professions 
    and states). Also follows the preposition \textit{z} (with) and \textit{za} 
    (behind/after in positional sense). Equivalent to Latin ablative of means.\\
    \textit{z \textbf{psem}} — with a dog\\
    \textit{z \textbf{wykształcenia} inżynier} — an engineer by training

  \item[Locative (\textit{miejscownik})] Location and topic. 
    \textbf{Always} used with a preposition (\textit{w, na, przy, o, po}).
    Never appears without one. Equivalent to Latin ablative of place.\\
    \textit{na \textbf{cmentarzu}} — at the cemetery\\
    \textit{przy \textbf{grobie}} — by the grave

  \item[Vocative (\textit{wołacz})] Direct address. Used when calling 
    or speaking to someone by name or title. Relatively rare in prose 
    but common in dialogue.\\
    \textit{Mamo!} — Mum! (addressing one's mother)

\end{description}

\grambox{Key principle}{Case endings change for both nouns \textit{and} 
adjectives — they must always agree in gender, number, and case. 
This is identical to Latin adjective-noun concord: 
\textit{wczesną jesienią} (in early autumn) — both words take the 
instrumental feminine singular ending.}

% ────────────────────────────────────────────────────────────
\section{Adjectives}

Adjectives agree with their nouns in gender, number, and case. The 
ending of the adjective therefore changes to match the noun it modifies, 
across all seven cases and three genders — 42 potential forms in principle, 
though many are identical in practice.

\textit{świeżo wykopany grób} (freshly dug grave, masc. nom.) — the participle 
\textit{wykopany} agrees with \textit{grób}.

\textit{świeżo wykopanym grobie} (by the freshly dug grave, masc. loc.) — both 
adjective and noun shift to the locative.

% ────────────────────────────────────────────────────────────
\section{Verbs: Aspect}

The single most important concept in Polish verbal grammar, and the one 
with no direct Latin equivalent, is \textbf{verbal aspect}.

Every Polish verb exists in two versions:

\begin{description}
  \item[Imperfective (\textit{niedokonany})] Describes ongoing, repeated, 
    or incomplete actions. What was happening, what happens habitually, 
    what is in progress.
  \item[Perfective (\textit{dokonany})] Describes completed, bounded, 
    or one-time actions. What happened (and is done), what will definitively happen.
\end{description}

\grambox{Key contrast}{\polword{pisać}{to write (imperfective, ongoing)} 
vs.\ \polword{napisać}{to write (perfective, completed)}. The prefix 
\textit{na-} here creates the perfective. This pattern — imperfective root, 
prefixed perfective — is one of the most common in Polish.}

Aspect affects \textbf{what tenses are possible}. Perfective verbs have 
no present tense (a completed action cannot be in progress now). Their 
``present-looking'' form is actually future: \textit{napiszę} = I will write (and finish it).

\subsection*{Common perfectivising prefixes}

These prefixes convert imperfective verbs to perfective and are extremely 
productive in Polish. Recognising them unlocks large swaths of vocabulary.

\begin{description}
  \item[\textbf{za-}] Often signals the onset or completion of an action: 
    \polword{pamiętać}{to remember} → \polword{zapomnieć}{to forget}; 
    \polword{truć}{to poison} → \polword{zatruć}{to poison (completely)}.
  \item[\textbf{po-}] Often signals completion or doing something for a while: 
    \polword{czekać}{to wait} → \polword{poczekać}{to wait (for a bit / done waiting)}.
  \item[\textbf{na-}] Completion, sufficiency: 
    \polword{pisać}{to write} → \polword{napisać}{to write (and finish)}.
  \item[\textbf{od-}] Away from, completing a departure: 
    \polword{iść}{to go} → \polword{odejść}{to go away, depart}.
  \item[\textbf{wy-}] Out, thoroughly: 
    \polword{kopać}{to dig} → \polword{wykopać}{to dig out/up}.
  \item[\textbf{z- / s-}] Completion, bringing together: 
    \polword{bliski}{close} → \polword{zbliżyć się}{to draw closer}.
\end{description}

% ────────────────────────────────────────────────────────────
\section{Verb Conjugation}

Polish verbs encode person and number in their endings, making subject 
pronouns optional (as in Latin and Russian). The pronouns are used 
for emphasis or contrast.

\subsection*{Present tense — example: \textit{czekać} (to wait, imperfective)}

\begin{tabular}{lll}
  \textbf{Person} & \textbf{Singular} & \textbf{Plural} \\
  1st & czekam (I wait) & czekamy (we wait) \\
  2nd & czekasz (you wait) & czekacie (you wait) \\
  3rd & czeka (he/she/it waits) & czekają (they wait) \\
\end{tabular}

\subsection*{Past tense — gender marking}

Uniquely among European languages, Polish past tense verbs agree with 
the gender of the subject. This gender marking is a suffix added to the 
past stem:

\begin{description}
  \item[Masculine singular] \textbf{-ł}: \textit{powiedział} (he said), 
    \textit{spotkałem} (I [masc.] met)
  \item[Feminine singular] \textbf{-ła}: \textit{czekała} (she waited), 
    \textit{musiała} (she had to)
  \item[Neuter singular] \textbf{-ło}: \textit{miało} (it was supposed to)
  \item[Virile plural (masc. personal)] \textbf{-li}: \textit{zdecydowaliśmy} 
    (we [men/mixed] decided), \textit{ustawili się} (they positioned themselves)
  \item[Non-virile plural] \textbf{-ły}: used for groups of women or non-persons
\end{description}

\grambox{Virile vs.\ non-virile}{Polish distinguishes between groups 
containing at least one man (\textit{virile}, personal masculine plural) 
and all other plurals. This distinction affects verb endings, adjective 
forms, and pronouns throughout the grammar.}

% ────────────────────────────────────────────────────────────
\section{Key Impersonal Constructions}

Polish uses several impersonal constructions that have no direct English 
equivalent but will be encountered on virtually every page.

\begin{description}
  \item[\polword{trzeba}{one must, it is necessary}] Followed by an 
    infinitive. No subject. \textit{Trzeba kontynuować} — one must continue. 
    Equivalent to French \textit{il faut}.
  \item[\polword{można}{one can, it is possible}] Followed by an 
    infinitive. \textit{Można było tam dostać} — one could get there. 
    Equivalent to French \textit{on peut}.
  \item[\polword{warto}{it is worth}] Followed by an infinitive. 
    \textit{Warto wiedzieć} — it is worth knowing.
\end{description}

% ────────────────────────────────────────────────────────────
\section{Essential Particles and Connectives}

A small set of particles carries enormous communicative weight in Polish 
and will be encountered constantly.

\begin{description}
  \item[\polword{już}{already, now, anymore}] Marks a transition or 
    change of state. \textit{Już rozumiem} — I understand now (I've crossed 
    into understanding). With negation: \textit{już nie} — not anymore.
  \item[\polword{jeszcze}{still, yet, even more}] Marks continuation 
    or addition. \textit{Jeszcze jeden} — yet one more. With negation: 
    \textit{jeszcze nie} — not yet.
  \item[\polword{już nie}{not anymore}] Something was happening but has stopped.
  \item[\polword{jeszcze nie}{not yet}] Something has not happened but may.
  \item[\polword{też / także}{also, too}] \textit{Też} is conversational; 
    \textit{także} is slightly more formal.
  \item[\polword{wcale nie}{not at all}] Emphatic negation. 
    \textit{Wcale nie odchodzi} — it does not go away at all.
  \item[\polword{właśnie}{exactly, just now, precisely}] 
    Confirms or emphasises: \textit{właśnie} as a standalone response 
    means ``exactly'' or ``quite so.''
  \item[\polword{chyba}{probably, I suppose, surely}] Expresses 
    uncertainty or mild assertion. Very common in speech.
  \item[\polword{przecież}{after all, but surely}] Implies that 
    something should be obvious. \textit{Przecież wiesz} — you know 
    this perfectly well (after all).
\end{description}

% ────────────────────────────────────────────────────────────
\section{Conjunctions: Because}

Three words translate as ``because'' with different registers:

\begin{description}
  \item[\polword{bo}{because}] Casual, spoken, everyday. 
    The default in conversation.
  \item[\polword{ponieważ}{because}] Standard written and formal spoken. 
    Neutral and widely applicable.
  \item[\polword{gdyż}{because/for}] Literary and elevated. 
    Signals deliberate, crafted prose. Equivalent to French \textit{car}.
\end{description}

% ────────────────────────────────────────────────────────────
\section{Word Order}

Polish word order is considerably freer than English, because case 
endings — not position — identify grammatical role. However, order 
is far from arbitrary: it encodes \textbf{information structure}.

\begin{description}
  \item[New information] tends to come at the \textbf{end} of the sentence 
    (the ``rheme'' or focus position). This is why Polish often places the 
    subject after the verb, or saves the most concrete noun for last.
  \item[Given / known information] tends to come early.
  \item[Verb-first] is used for dramatic effect, as in: 
    \textit{Umarł Fiszer} — Fiszer died (the death before the name).
\end{description}

This freedom means that close translation often fails to capture 
the author's emphasis. Notes in the chapter analyses point out 
significant word-order choices where they carry meaning.

% ────────────────────────────────────────────────────────────
\section{Reflexive Verbs and \textit{się}}

The particle \textbf{się} is one of the most versatile elements in 
Polish. It functions as:

\begin{description}
  \item[Reflexive marker] The action reflects back on the subject: 
    \textit{zbliżyć się} — to bring oneself closer, to approach.
  \item[Reciprocal marker] \textit{spotkać się} — to meet (each other).
  \item[Impersonal/passive marker] \textit{mówi się} — one says, 
    it is said. Removes the agent entirely.
  \item[Part of the verb's inherent meaning] Some verbs simply 
    require \textit{się} with no separable reflexive meaning: 
    \textit{bać się} (to be afraid), \textit{cieszyć się} (to be glad).
\end{description}

% ────────────────────────────────────────────────────────────
\section{Participles}

Polish uses participial constructions extensively — more so than 
conversational English, but comparable to written French or Latin.

\begin{description}
  \item[Present active participle (\textit{imiesłów przymiotnikowy czynny})] 
    Describes an ongoing action of the noun it modifies. 
    Formed with suffix \textbf{-ący/-ąca/-ące}: 
    \textit{grupka \textbf{stojąca} przy grobie} — the group standing by the grave.
  \item[Past passive participle (\textit{imiesłów przymiotnikowy bierny})] 
    Describes a completed action done to the noun. 
    Formed with suffix \textbf{-ny/-na/-ne} or \textbf{-ty/-ta/-te}: 
    \textit{świeżo \textbf{wykopany} grób} — the freshly dug grave; 
    \textit{zapomniana epidemia} — the forgotten epidemic.
  \item[Adverbial participle (\textit{imiesłów przysłówkowy})] 
    Describes a simultaneous action of the subject. Does not decline. 
    Formed with suffix \textbf{-ąc}: 
    \textit{\textbf{odkręcając} nakrętkę} — while unscrewing the cap.
\end{description}

% ============================================================

\chapter*{Polish Grammar Foundations}
\addcontentsline{toc}{chapter}{Polish Grammar Foundations}
\markboth{Polish Grammar Foundations}{}

This chapter provides a compact reference for the grammatical structures 
encountered throughout the reading. It is not intended as a complete grammar 
— rather, as a touchstone for constructs that recur in the text. Cross-references 
to passages where each construct first appears are noted where relevant.

% ────────────────────────────────────────────────────────────
\section{Nouns and Gender}

Polish nouns have three grammatical genders: \textbf{masculine}, \textbf{feminine}, 
and \textbf{neuter}. Unlike French, where gender must largely be memorised, Polish 
gender is often inferable from the noun's ending.

\begin{description}
  \item[Masculine] Typically end in a consonant: \polword{grób}{grave}, 
    \polword{pogrzeb}{funeral}, \polword{syn}{son}.
  \item[Feminine] Typically end in \textbf{-a} or \textbf{-i/y}: 
    \polword{wdowa}{widow}, \polword{grupka}{small group}, 
    \polword{nadzieja}{hope}.
  \item[Neuter] Typically end in \textbf{-o} or \textbf{-e}: 
    \polword{dzieciństwo}{childhood}, \polword{dziecko}{child}.
\end{description}

\grambox{Note on animacy}{Masculine nouns are further divided into 
\textit{animate} (persons and animals) and \textit{inanimate}. This distinction 
affects the accusative case ending and is particularly important for direct objects.}

Nouns also form diminutives extensively in Polish, often with suffixes 
\textbf{-ka}, \textbf{-ko}, \textbf{-ek}, \textbf{-eczka}. These indicate 
smallness but frequently carry emotional colouring — affection, sympathy, or 
irony — rather than literal size. Examples encountered in Chapter 1: 
\polword{grupka}{little group} (from \textit{grupa}), 
\polword{garstka}{small handful} (from \textit{garść}), 
\polword{flaszka}{bottle, colloquial} (from \textit{flasza}),
\polword{podstawówka}{primary school, affectionate} (from \textit{szkoła podstawowa}).

% ────────────────────────────────────────────────────────────
\section{The Seven Cases}

Polish has seven grammatical cases. Each case encodes a grammatical 
relationship that English expresses through word order or prepositions. 
The Latin equivalents are noted as a reference point.

\begin{description}

  \item[Nominative (\textit{mianownik})] The subject of the sentence. 
    Who or what performs the action. Equivalent to Latin nominative.\\
    \textit{Groszek \textbf{powiedział}\ldots} — Groszek said\ldots

  \item[Genitive (\textit{dopełniacz})] Possession, absence, quantity, 
    and the object of many prepositions (\textit{do, od, z, bez, dla, oprócz}). 
    The most frequently used case after the nominative. Equivalent to Latin genitive.\\
    \textit{dla \textbf{uczczenia} jego \textbf{pamięci}} — for the honouring of his memory\\
    \textit{garstki \textbf{obecnych}} — a handful of those present

  \item[Dative (\textit{celownik})] The indirect object — to whom or for whom. 
    Equivalent to Latin dative. Encodes dedication, as seen in the book's 
    epigraph: \textit{\textbf{Ojcu}, \textbf{Synowi}} — to [my] father, to [my] son.

  \item[Accusative (\textit{biernik})] The direct object of most verbs, 
    and the object of motion prepositions (\textit{na, w, przez, za}).
    Equivalent to Latin accusative.\\
    \textit{Zostaw \textbf{tę flaszkę}!} — Leave that bottle!

  \item[Instrumental (\textit{narzędnik})] Means, accompaniment, and 
    identity (used after the verb \textit{być} — to be, for professions 
    and states). Also follows the preposition \textit{z} (with) and \textit{za} 
    (behind/after in positional sense). Equivalent to Latin ablative of means.\\
    \textit{z \textbf{psem}} — with a dog\\
    \textit{z \textbf{wykształcenia} inżynier} — an engineer by training

  \item[Locative (\textit{miejscownik})] Location and topic. 
    \textbf{Always} used with a preposition (\textit{w, na, przy, o, po}).
    Never appears without one. Equivalent to Latin ablative of place.\\
    \textit{na \textbf{cmentarzu}} — at the cemetery\\
    \textit{przy \textbf{grobie}} — by the grave

  \item[Vocative (\textit{wołacz})] Direct address. Used when calling 
    or speaking to someone by name or title. Relatively rare in prose 
    but common in dialogue.\\
    \textit{Mamo!} — Mum! (addressing one's mother)

\end{description}

\grambox{Key principle}{Case endings change for both nouns \textit{and} 
adjectives — they must always agree in gender, number, and case. 
This is identical to Latin adjective-noun concord: 
\textit{wczesną jesienią} (in early autumn) — both words take the 
instrumental feminine singular ending.}

% ────────────────────────────────────────────────────────────
\section{Adjectives}

Adjectives agree with their nouns in gender, number, and case. The 
ending of the adjective therefore changes to match the noun it modifies, 
across all seven cases and three genders — 42 potential forms in principle, 
though many are identical in practice.

\textit{świeżo wykopany grób} (freshly dug grave, masc. nom.) — the participle 
\textit{wykopany} agrees with \textit{grób}.

\textit{świeżo wykopanym grobie} (by the freshly dug grave, masc. loc.) — both 
adjective and noun shift to the locative.

% ────────────────────────────────────────────────────────────
\section{Verbs: Aspect}

The single most important concept in Polish verbal grammar, and the one 
with no direct Latin equivalent, is \textbf{verbal aspect}.

Every Polish verb exists in two versions:

\begin{description}
  \item[Imperfective (\textit{niedokonany})] Describes ongoing, repeated, 
    or incomplete actions. What was happening, what happens habitually, 
    what is in progress.
  \item[Perfective (\textit{dokonany})] Describes completed, bounded, 
    or one-time actions. What happened (and is done), what will definitively happen.
\end{description}

\grambox{Key contrast}{\polword{pisać}{to write (imperfective, ongoing)} 
vs.\ \polword{napisać}{to write (perfective, completed)}. The prefix 
\textit{na-} here creates the perfective. This pattern — imperfective root, 
prefixed perfective — is one of the most common in Polish.}

Aspect affects \textbf{what tenses are possible}. Perfective verbs have 
no present tense (a completed action cannot be in progress now). Their 
``present-looking'' form is actually future: \textit{napiszę} = I will write (and finish it).

\subsection*{Common perfectivising prefixes}

These prefixes convert imperfective verbs to perfective and are extremely 
productive in Polish. Recognising them unlocks large swaths of vocabulary.

\begin{description}
  \item[\textbf{za-}] Often signals the onset or completion of an action: 
    \polword{pamiętać}{to remember} → \polword{zapomnieć}{to forget}; 
    \polword{truć}{to poison} → \polword{zatruć}{to poison (completely)}.
  \item[\textbf{po-}] Often signals completion or doing something for a while: 
    \polword{czekać}{to wait} → \polword{poczekać}{to wait (for a bit / done waiting)}.
  \item[\textbf{na-}] Completion, sufficiency: 
    \polword{pisać}{to write} → \polword{napisać}{to write (and finish)}.
  \item[\textbf{od-}] Away from, completing a departure: 
    \polword{iść}{to go} → \polword{odejść}{to go away, depart}.
  \item[\textbf{wy-}] Out, thoroughly: 
    \polword{kopać}{to dig} → \polword{wykopać}{to dig out/up}.
  \item[\textbf{z- / s-}] Completion, bringing together: 
    \polword{bliski}{close} → \polword{zbliżyć się}{to draw closer}.
\end{description}

% ────────────────────────────────────────────────────────────
\section{Verb Conjugation}

Polish verbs encode person and number in their endings, making subject 
pronouns optional (as in Latin and Russian). The pronouns are used 
for emphasis or contrast.

\subsection*{Present tense — example: \textit{czekać} (to wait, imperfective)}

\begin{tabular}{lll}
  \textbf{Person} & \textbf{Singular} & \textbf{Plural} \\
  1st & czekam (I wait) & czekamy (we wait) \\
  2nd & czekasz (you wait) & czekacie (you wait) \\
  3rd & czeka (he/she/it waits) & czekają (they wait) \\
\end{tabular}

\subsection*{Past tense — gender marking}

Uniquely among European languages, Polish past tense verbs agree with 
the gender of the subject. This gender marking is a suffix added to the 
past stem:

\begin{description}
  \item[Masculine singular] \textbf{-ł}: \textit{powiedział} (he said), 
    \textit{spotkałem} (I [masc.] met)
  \item[Feminine singular] \textbf{-ła}: \textit{czekała} (she waited), 
    \textit{musiała} (she had to)
  \item[Neuter singular] \textbf{-ło}: \textit{miało} (it was supposed to)
  \item[Virile plural (masc. personal)] \textbf{-li}: \textit{zdecydowaliśmy} 
    (we [men/mixed] decided), \textit{ustawili się} (they positioned themselves)
  \item[Non-virile plural] \textbf{-ły}: used for groups of women or non-persons
\end{description}

\grambox{Virile vs.\ non-virile}{Polish distinguishes between groups 
containing at least one man (\textit{virile}, personal masculine plural) 
and all other plurals. This distinction affects verb endings, adjective 
forms, and pronouns throughout the grammar.}

% ────────────────────────────────────────────────────────────
\section{Key Impersonal Constructions}

Polish uses several impersonal constructions that have no direct English 
equivalent but will be encountered on virtually every page.

\begin{description}
  \item[\polword{trzeba}{one must, it is necessary}] Followed by an 
    infinitive. No subject. \textit{Trzeba kontynuować} — one must continue. 
    Equivalent to French \textit{il faut}.
  \item[\polword{można}{one can, it is possible}] Followed by an 
    infinitive. \textit{Można było tam dostać} — one could get there. 
    Equivalent to French \textit{on peut}.
  \item[\polword{warto}{it is worth}] Followed by an infinitive. 
    \textit{Warto wiedzieć} — it is worth knowing.
\end{description}

% ────────────────────────────────────────────────────────────
\section{Essential Particles and Connectives}

A small set of particles carries enormous communicative weight in Polish 
and will be encountered constantly.

\begin{description}
  \item[\polword{już}{already, now, anymore}] Marks a transition or 
    change of state. \textit{Już rozumiem} — I understand now (I've crossed 
    into understanding). With negation: \textit{już nie} — not anymore.
  \item[\polword{jeszcze}{still, yet, even more}] Marks continuation 
    or addition. \textit{Jeszcze jeden} — yet one more. With negation: 
    \textit{jeszcze nie} — not yet.
  \item[\polword{już nie}{not anymore}] Something was happening but has stopped.
  \item[\polword{jeszcze nie}{not yet}] Something has not happened but may.
  \item[\polword{też / także}{also, too}] \textit{Też} is conversational; 
    \textit{także} is slightly more formal.
  \item[\polword{wcale nie}{not at all}] Emphatic negation. 
    \textit{Wcale nie odchodzi} — it does not go away at all.
  \item[\polword{właśnie}{exactly, just now, precisely}] 
    Confirms or emphasises: \textit{właśnie} as a standalone response 
    means ``exactly'' or ``quite so.''
  \item[\polword{chyba}{probably, I suppose, surely}] Expresses 
    uncertainty or mild assertion. Very common in speech.
  \item[\polword{przecież}{after all, but surely}] Implies that 
    something should be obvious. \textit{Przecież wiesz} — you know 
    this perfectly well (after all).
\end{description}

% ────────────────────────────────────────────────────────────
\section{Conjunctions: Because}

Three words translate as ``because'' with different registers:

\begin{description}
  \item[\polword{bo}{because}] Casual, spoken, everyday. 
    The default in conversation.
  \item[\polword{ponieważ}{because}] Standard written and formal spoken. 
    Neutral and widely applicable.
  \item[\polword{gdyż}{because/for}] Literary and elevated. 
    Signals deliberate, crafted prose. Equivalent to French \textit{car}.
\end{description}

% ────────────────────────────────────────────────────────────
\section{Word Order}

Polish word order is considerably freer than English, because case 
endings — not position — identify grammatical role. However, order 
is far from arbitrary: it encodes \textbf{information structure}.

\begin{description}
  \item[New information] tends to come at the \textbf{end} of the sentence 
    (the ``rheme'' or focus position). This is why Polish often places the 
    subject after the verb, or saves the most concrete noun for last.
  \item[Given / known information] tends to come early.
  \item[Verb-first] is used for dramatic effect, as in: 
    \textit{Umarł Fiszer} — Fiszer died (the death before the name).
\end{description}

This freedom means that close translation often fails to capture 
the author's emphasis. Notes in the chapter analyses point out 
significant word-order choices where they carry meaning.

% ────────────────────────────────────────────────────────────
\section{Reflexive Verbs and \textit{się}}

The particle \textbf{się} is one of the most versatile elements in 
Polish. It functions as:

\begin{description}
  \item[Reflexive marker] The action reflects back on the subject: 
    \textit{zbliżyć się} — to bring oneself closer, to approach.
  \item[Reciprocal marker] \textit{spotkać się} — to meet (each other).
  \item[Impersonal/passive marker] \textit{mówi się} — one says, 
    it is said. Removes the agent entirely.
  \item[Part of the verb's inherent meaning] Some verbs simply 
    require \textit{się} with no separable reflexive meaning: 
    \textit{bać się} (to be afraid), \textit{cieszyć się} (to be glad).
\end{description}

% ────────────────────────────────────────────────────────────
\section{Participles}

Polish uses participial constructions extensively — more so than 
conversational English, but comparable to written French or Latin.

\begin{description}
  \item[Present active participle (\textit{imiesłów przymiotnikowy czynny})] 
    Describes an ongoing action of the noun it modifies. 
    Formed with suffix \textbf{-ący/-ąca/-ące}: 
    \textit{grupka \textbf{stojąca} przy grobie} — the group standing by the grave.
  \item[Past passive participle (\textit{imiesłów przymiotnikowy bierny})] 
    Describes a completed action done to the noun. 
    Formed with suffix \textbf{-ny/-na/-ne} or \textbf{-ty/-ta/-te}: 
    \textit{świeżo \textbf{wykopany} grób} — the freshly dug grave; 
    \textit{zapomniana epidemia} — the forgotten epidemic.
  \item[Adverbial participle (\textit{imiesłów przysłówkowy})] 
    Describes a simultaneous action of the subject. Does not decline. 
    Formed with suffix \textbf{-ąc}: 
    \textit{\textbf{odkręcając} nakrętkę} — while unscrewing the cap.
\end{description}

% ============================================================

\chapter*{Polish Grammar Foundations}
\addcontentsline{toc}{chapter}{Polish Grammar Foundations}
\markboth{Polish Grammar Foundations}{}

This chapter provides a compact reference for the grammatical structures 
encountered throughout the reading. It is not intended as a complete grammar 
— rather, as a touchstone for constructs that recur in the text. Cross-references 
to passages where each construct first appears are noted where relevant.

% ────────────────────────────────────────────────────────────
\section{Nouns and Gender}

Polish nouns have three grammatical genders: \textbf{masculine}, \textbf{feminine}, 
and \textbf{neuter}. Unlike French, where gender must largely be memorised, Polish 
gender is often inferable from the noun's ending.

\begin{description}
  \item[Masculine] Typically end in a consonant: \polword{grób}{grave}, 
    \polword{pogrzeb}{funeral}, \polword{syn}{son}.
  \item[Feminine] Typically end in \textbf{-a} or \textbf{-i/y}: 
    \polword{wdowa}{widow}, \polword{grupka}{small group}, 
    \polword{nadzieja}{hope}.
  \item[Neuter] Typically end in \textbf{-o} or \textbf{-e}: 
    \polword{dzieciństwo}{childhood}, \polword{dziecko}{child}.
\end{description}

\grambox{Note on animacy}{Masculine nouns are further divided into 
\textit{animate} (persons and animals) and \textit{inanimate}. This distinction 
affects the accusative case ending and is particularly important for direct objects.}

Nouns also form diminutives extensively in Polish, often with suffixes 
\textbf{-ka}, \textbf{-ko}, \textbf{-ek}, \textbf{-eczka}. These indicate 
smallness but frequently carry emotional colouring — affection, sympathy, or 
irony — rather than literal size. Examples encountered in Chapter 1: 
\polword{grupka}{little group} (from \textit{grupa}), 
\polword{garstka}{small handful} (from \textit{garść}), 
\polword{flaszka}{bottle, colloquial} (from \textit{flasza}),
\polword{podstawówka}{primary school, affectionate} (from \textit{szkoła podstawowa}).

% ────────────────────────────────────────────────────────────
\section{The Seven Cases}

Polish has seven grammatical cases. Each case encodes a grammatical 
relationship that English expresses through word order or prepositions. 
The Latin equivalents are noted as a reference point.

\begin{description}

  \item[Nominative (\textit{mianownik})] The subject of the sentence. 
    Who or what performs the action. Equivalent to Latin nominative.\\
    \textit{Groszek \textbf{powiedział}\ldots} — Groszek said\ldots

  \item[Genitive (\textit{dopełniacz})] Possession, absence, quantity, 
    and the object of many prepositions (\textit{do, od, z, bez, dla, oprócz}). 
    The most frequently used case after the nominative. Equivalent to Latin genitive.\\
    \textit{dla \textbf{uczczenia} jego \textbf{pamięci}} — for the honouring of his memory\\
    \textit{garstki \textbf{obecnych}} — a handful of those present

  \item[Dative (\textit{celownik})] The indirect object — to whom or for whom. 
    Equivalent to Latin dative. Encodes dedication, as seen in the book's 
    epigraph: \textit{\textbf{Ojcu}, \textbf{Synowi}} — to [my] father, to [my] son.

  \item[Accusative (\textit{biernik})] The direct object of most verbs, 
    and the object of motion prepositions (\textit{na, w, przez, za}).
    Equivalent to Latin accusative.\\
    \textit{Zostaw \textbf{tę flaszkę}!} — Leave that bottle!

  \item[Instrumental (\textit{narzędnik})] Means, accompaniment, and 
    identity (used after the verb \textit{być} — to be, for professions 
    and states). Also follows the preposition \textit{z} (with) and \textit{za} 
    (behind/after in positional sense). Equivalent to Latin ablative of means.\\
    \textit{z \textbf{psem}} — with a dog\\
    \textit{z \textbf{wykształcenia} inżynier} — an engineer by training

  \item[Locative (\textit{miejscownik})] Location and topic. 
    \textbf{Always} used with a preposition (\textit{w, na, przy, o, po}).
    Never appears without one. Equivalent to Latin ablative of place.\\
    \textit{na \textbf{cmentarzu}} — at the cemetery\\
    \textit{przy \textbf{grobie}} — by the grave

  \item[Vocative (\textit{wołacz})] Direct address. Used when calling 
    or speaking to someone by name or title. Relatively rare in prose 
    but common in dialogue.\\
    \textit{Mamo!} — Mum! (addressing one's mother)

\end{description}

\grambox{Key principle}{Case endings change for both nouns \textit{and} 
adjectives — they must always agree in gender, number, and case. 
This is identical to Latin adjective-noun concord: 
\textit{wczesną jesienią} (in early autumn) — both words take the 
instrumental feminine singular ending.}

% ────────────────────────────────────────────────────────────
\section{Adjectives}

Adjectives agree with their nouns in gender, number, and case. The 
ending of the adjective therefore changes to match the noun it modifies, 
across all seven cases and three genders — 42 potential forms in principle, 
though many are identical in practice.

\textit{świeżo wykopany grób} (freshly dug grave, masc. nom.) — the participle 
\textit{wykopany} agrees with \textit{grób}.

\textit{świeżo wykopanym grobie} (by the freshly dug grave, masc. loc.) — both 
adjective and noun shift to the locative.

% ────────────────────────────────────────────────────────────
\section{Verbs: Aspect}

The single most important concept in Polish verbal grammar, and the one 
with no direct Latin equivalent, is \textbf{verbal aspect}.

Every Polish verb exists in two versions:

\begin{description}
  \item[Imperfective (\textit{niedokonany})] Describes ongoing, repeated, 
    or incomplete actions. What was happening, what happens habitually, 
    what is in progress.
  \item[Perfective (\textit{dokonany})] Describes completed, bounded, 
    or one-time actions. What happened (and is done), what will definitively happen.
\end{description}

\grambox{Key contrast}{\polword{pisać}{to write (imperfective, ongoing)} 
vs.\ \polword{napisać}{to write (perfective, completed)}. The prefix 
\textit{na-} here creates the perfective. This pattern — imperfective root, 
prefixed perfective — is one of the most common in Polish.}

Aspect affects \textbf{what tenses are possible}. Perfective verbs have 
no present tense (a completed action cannot be in progress now). Their 
``present-looking'' form is actually future: \textit{napiszę} = I will write (and finish it).

\subsection*{Common perfectivising prefixes}

These prefixes convert imperfective verbs to perfective and are extremely 
productive in Polish. Recognising them unlocks large swaths of vocabulary.

\begin{description}
  \item[\textbf{za-}] Often signals the onset or completion of an action: 
    \polword{pamiętać}{to remember} → \polword{zapomnieć}{to forget}; 
    \polword{truć}{to poison} → \polword{zatruć}{to poison (completely)}.
  \item[\textbf{po-}] Often signals completion or doing something for a while: 
    \polword{czekać}{to wait} → \polword{poczekać}{to wait (for a bit / done waiting)}.
  \item[\textbf{na-}] Completion, sufficiency: 
    \polword{pisać}{to write} → \polword{napisać}{to write (and finish)}.
  \item[\textbf{od-}] Away from, completing a departure: 
    \polword{iść}{to go} → \polword{odejść}{to go away, depart}.
  \item[\textbf{wy-}] Out, thoroughly: 
    \polword{kopać}{to dig} → \polword{wykopać}{to dig out/up}.
  \item[\textbf{z- / s-}] Completion, bringing together: 
    \polword{bliski}{close} → \polword{zbliżyć się}{to draw closer}.
\end{description}

% ────────────────────────────────────────────────────────────
\section{Verb Conjugation}

Polish verbs encode person and number in their endings, making subject 
pronouns optional (as in Latin and Russian). The pronouns are used 
for emphasis or contrast.

\subsection*{Present tense — example: \textit{czekać} (to wait, imperfective)}

\begin{tabular}{lll}
  \textbf{Person} & \textbf{Singular} & \textbf{Plural} \\
  1st & czekam (I wait) & czekamy (we wait) \\
  2nd & czekasz (you wait) & czekacie (you wait) \\
  3rd & czeka (he/she/it waits) & czekają (they wait) \\
\end{tabular}

\subsection*{Past tense — gender marking}

Uniquely among European languages, Polish past tense verbs agree with 
the gender of the subject. This gender marking is a suffix added to the 
past stem:

\begin{description}
  \item[Masculine singular] \textbf{-ł}: \textit{powiedział} (he said), 
    \textit{spotkałem} (I [masc.] met)
  \item[Feminine singular] \textbf{-ła}: \textit{czekała} (she waited), 
    \textit{musiała} (she had to)
  \item[Neuter singular] \textbf{-ło}: \textit{miało} (it was supposed to)
  \item[Virile plural (masc. personal)] \textbf{-li}: \textit{zdecydowaliśmy} 
    (we [men/mixed] decided), \textit{ustawili się} (they positioned themselves)
  \item[Non-virile plural] \textbf{-ły}: used for groups of women or non-persons
\end{description}

\grambox{Virile vs.\ non-virile}{Polish distinguishes between groups 
containing at least one man (\textit{virile}, personal masculine plural) 
and all other plurals. This distinction affects verb endings, adjective 
forms, and pronouns throughout the grammar.}

% ────────────────────────────────────────────────────────────
\section{Key Impersonal Constructions}

Polish uses several impersonal constructions that have no direct English 
equivalent but will be encountered on virtually every page.

\begin{description}
  \item[\polword{trzeba}{one must, it is necessary}] Followed by an 
    infinitive. No subject. \textit{Trzeba kontynuować} — one must continue. 
    Equivalent to French \textit{il faut}.
  \item[\polword{można}{one can, it is possible}] Followed by an 
    infinitive. \textit{Można było tam dostać} — one could get there. 
    Equivalent to French \textit{on peut}.
  \item[\polword{warto}{it is worth}] Followed by an infinitive. 
    \textit{Warto wiedzieć} — it is worth knowing.
\end{description}

% ────────────────────────────────────────────────────────────
\section{Essential Particles and Connectives}

A small set of particles carries enormous communicative weight in Polish 
and will be encountered constantly.

\begin{description}
  \item[\polword{już}{already, now, anymore}] Marks a transition or 
    change of state. \textit{Już rozumiem} — I understand now (I've crossed 
    into understanding). With negation: \textit{już nie} — not anymore.
  \item[\polword{jeszcze}{still, yet, even more}] Marks continuation 
    or addition. \textit{Jeszcze jeden} — yet one more. With negation: 
    \textit{jeszcze nie} — not yet.
  \item[\polword{już nie}{not anymore}] Something was happening but has stopped.
  \item[\polword{jeszcze nie}{not yet}] Something has not happened but may.
  \item[\polword{też / także}{also, too}] \textit{Też} is conversational; 
    \textit{także} is slightly more formal.
  \item[\polword{wcale nie}{not at all}] Emphatic negation. 
    \textit{Wcale nie odchodzi} — it does not go away at all.
  \item[\polword{właśnie}{exactly, just now, precisely}] 
    Confirms or emphasises: \textit{właśnie} as a standalone response 
    means ``exactly'' or ``quite so.''
  \item[\polword{chyba}{probably, I suppose, surely}] Expresses 
    uncertainty or mild assertion. Very common in speech.
  \item[\polword{przecież}{after all, but surely}] Implies that 
    something should be obvious. \textit{Przecież wiesz} — you know 
    this perfectly well (after all).
\end{description}

% ────────────────────────────────────────────────────────────
\section{Conjunctions: Because}

Three words translate as ``because'' with different registers:

\begin{description}
  \item[\polword{bo}{because}] Casual, spoken, everyday. 
    The default in conversation.
  \item[\polword{ponieważ}{because}] Standard written and formal spoken. 
    Neutral and widely applicable.
  \item[\polword{gdyż}{because/for}] Literary and elevated. 
    Signals deliberate, crafted prose. Equivalent to French \textit{car}.
\end{description}

% ────────────────────────────────────────────────────────────
\section{Word Order}

Polish word order is considerably freer than English, because case 
endings — not position — identify grammatical role. However, order 
is far from arbitrary: it encodes \textbf{information structure}.

\begin{description}
  \item[New information] tends to come at the \textbf{end} of the sentence 
    (the ``rheme'' or focus position). This is why Polish often places the 
    subject after the verb, or saves the most concrete noun for last.
  \item[Given / known information] tends to come early.
  \item[Verb-first] is used for dramatic effect, as in: 
    \textit{Umarł Fiszer} — Fiszer died (the death before the name).
\end{description}

This freedom means that close translation often fails to capture 
the author's emphasis. Notes in the chapter analyses point out 
significant word-order choices where they carry meaning.

% ────────────────────────────────────────────────────────────
\section{Reflexive Verbs and \textit{się}}

The particle \textbf{się} is one of the most versatile elements in 
Polish. It functions as:

\begin{description}
  \item[Reflexive marker] The action reflects back on the subject: 
    \textit{zbliżyć się} — to bring oneself closer, to approach.
  \item[Reciprocal marker] \textit{spotkać się} — to meet (each other).
  \item[Impersonal/passive marker] \textit{mówi się} — one says, 
    it is said. Removes the agent entirely.
  \item[Part of the verb's inherent meaning] Some verbs simply 
    require \textit{się} with no separable reflexive meaning: 
    \textit{bać się} (to be afraid), \textit{cieszyć się} (to be glad).
\end{description}

% ────────────────────────────────────────────────────────────
\section{Participles}

Polish uses participial constructions extensively — more so than 
conversational English, but comparable to written French or Latin.

\begin{description}
  \item[Present active participle (\textit{imiesłów przymiotnikowy czynny})] 
    Describes an ongoing action of the noun it modifies. 
    Formed with suffix \textbf{-ący/-ąca/-ące}: 
    \textit{grupka \textbf{stojąca} przy grobie} — the group standing by the grave.
  \item[Past passive participle (\textit{imiesłów przymiotnikowy bierny})] 
    Describes a completed action done to the noun. 
    Formed with suffix \textbf{-ny/-na/-ne} or \textbf{-ty/-ta/-te}: 
    \textit{świeżo \textbf{wykopany} grób} — the freshly dug grave; 
    \textit{zapomniana epidemia} — the forgotten epidemic.
  \item[Adverbial participle (\textit{imiesłów przysłówkowy})] 
    Describes a simultaneous action of the subject. Does not decline. 
    Formed with suffix \textbf{-ąc}: 
    \textit{\textbf{odkręcając} nakrętkę} — while unscrewing the cap.
\end{description}

% ============================================================

\chapter*{Polish Grammar Foundations}
\addcontentsline{toc}{chapter}{Polish Grammar Foundations}
\markboth{Polish Grammar Foundations}{}

This chapter provides a compact reference for the grammatical structures 
encountered throughout the reading. It is not intended as a complete grammar 
— rather, as a touchstone for constructs that recur in the text. Cross-references 
to passages where each construct first appears are noted where relevant.

% ────────────────────────────────────────────────────────────
\section{Nouns and Gender}

Polish nouns have three grammatical genders: \textbf{masculine}, \textbf{feminine}, 
and \textbf{neuter}. Unlike French, where gender must largely be memorised, Polish 
gender is often inferable from the noun's ending.

\begin{description}
  \item[Masculine] Typically end in a consonant: \polword{grób}{grave}, 
    \polword{pogrzeb}{funeral}, \polword{syn}{son}.
  \item[Feminine] Typically end in \textbf{-a} or \textbf{-i/y}: 
    \polword{wdowa}{widow}, \polword{grupka}{small group}, 
    \polword{nadzieja}{hope}.
  \item[Neuter] Typically end in \textbf{-o} or \textbf{-e}: 
    \polword{dzieciństwo}{childhood}, \polword{dziecko}{child}.
\end{description}

\grambox{Note on animacy}{Masculine nouns are further divided into 
\textit{animate} (persons and animals) and \textit{inanimate}. This distinction 
affects the accusative case ending and is particularly important for direct objects.}

Nouns also form diminutives extensively in Polish, often with suffixes 
\textbf{-ka}, \textbf{-ko}, \textbf{-ek}, \textbf{-eczka}. These indicate 
smallness but frequently carry emotional colouring — affection, sympathy, or 
irony — rather than literal size. Examples encountered in Chapter 1: 
\polword{grupka}{little group} (from \textit{grupa}), 
\polword{garstka}{small handful} (from \textit{garść}), 
\polword{flaszka}{bottle, colloquial} (from \textit{flasza}),
\polword{podstawówka}{primary school, affectionate} (from \textit{szkoła podstawowa}).

% ────────────────────────────────────────────────────────────
\section{The Seven Cases}

Polish has seven grammatical cases. Each case encodes a grammatical 
relationship that English expresses through word order or prepositions. 
The Latin equivalents are noted as a reference point.

\begin{description}

  \item[Nominative (\textit{mianownik})] The subject of the sentence. 
    Who or what performs the action. Equivalent to Latin nominative.\\
    \textit{Groszek \textbf{powiedział}\ldots} — Groszek said\ldots

  \item[Genitive (\textit{dopełniacz})] Possession, absence, quantity, 
    and the object of many prepositions (\textit{do, od, z, bez, dla, oprócz}). 
    The most frequently used case after the nominative. Equivalent to Latin genitive.\\
    \textit{dla \textbf{uczczenia} jego \textbf{pamięci}} — for the honouring of his memory\\
    \textit{garstki \textbf{obecnych}} — a handful of those present

  \item[Dative (\textit{celownik})] The indirect object — to whom or for whom. 
    Equivalent to Latin dative. Encodes dedication, as seen in the book's 
    epigraph: \textit{\textbf{Ojcu}, \textbf{Synowi}} — to [my] father, to [my] son.

  \item[Accusative (\textit{biernik})] The direct object of most verbs, 
    and the object of motion prepositions (\textit{na, w, przez, za}).
    Equivalent to Latin accusative.\\
    \textit{Zostaw \textbf{tę flaszkę}!} — Leave that bottle!

  \item[Instrumental (\textit{narzędnik})] Means, accompaniment, and 
    identity (used after the verb \textit{być} — to be, for professions 
    and states). Also follows the preposition \textit{z} (with) and \textit{za} 
    (behind/after in positional sense). Equivalent to Latin ablative of means.\\
    \textit{z \textbf{psem}} — with a dog\\
    \textit{z \textbf{wykształcenia} inżynier} — an engineer by training

  \item[Locative (\textit{miejscownik})] Location and topic. 
    \textbf{Always} used with a preposition (\textit{w, na, przy, o, po}).
    Never appears without one. Equivalent to Latin ablative of place.\\
    \textit{na \textbf{cmentarzu}} — at the cemetery\\
    \textit{przy \textbf{grobie}} — by the grave

  \item[Vocative (\textit{wołacz})] Direct address. Used when calling 
    or speaking to someone by name or title. Relatively rare in prose 
    but common in dialogue.\\
    \textit{Mamo!} — Mum! (addressing one's mother)

\end{description}

\grambox{Key principle}{Case endings change for both nouns \textit{and} 
adjectives — they must always agree in gender, number, and case. 
This is identical to Latin adjective-noun concord: 
\textit{wczesną jesienią} (in early autumn) — both words take the 
instrumental feminine singular ending.}

% ────────────────────────────────────────────────────────────
\section{Adjectives}

Adjectives agree with their nouns in gender, number, and case. The 
ending of the adjective therefore changes to match the noun it modifies, 
across all seven cases and three genders — 42 potential forms in principle, 
though many are identical in practice.

\textit{świeżo wykopany grób} (freshly dug grave, masc. nom.) — the participle 
\textit{wykopany} agrees with \textit{grób}.

\textit{świeżo wykopanym grobie} (by the freshly dug grave, masc. loc.) — both 
adjective and noun shift to the locative.

% ────────────────────────────────────────────────────────────
\section{Verbs: Aspect}

The single most important concept in Polish verbal grammar, and the one 
with no direct Latin equivalent, is \textbf{verbal aspect}.

Every Polish verb exists in two versions:

\begin{description}
  \item[Imperfective (\textit{niedokonany})] Describes ongoing, repeated, 
    or incomplete actions. What was happening, what happens habitually, 
    what is in progress.
  \item[Perfective (\textit{dokonany})] Describes completed, bounded, 
    or one-time actions. What happened (and is done), what will definitively happen.
\end{description}

\grambox{Key contrast}{\polword{pisać}{to write (imperfective, ongoing)} 
vs.\ \polword{napisać}{to write (perfective, completed)}. The prefix 
\textit{na-} here creates the perfective. This pattern — imperfective root, 
prefixed perfective — is one of the most common in Polish.}

Aspect affects \textbf{what tenses are possible}. Perfective verbs have 
no present tense (a completed action cannot be in progress now). Their 
``present-looking'' form is actually future: \textit{napiszę} = I will write (and finish it).

\subsection*{Common perfectivising prefixes}

These prefixes convert imperfective verbs to perfective and are extremely 
productive in Polish. Recognising them unlocks large swaths of vocabulary.

\begin{description}
  \item[\textbf{za-}] Often signals the onset or completion of an action: 
    \polword{pamiętać}{to remember} → \polword{zapomnieć}{to forget}; 
    \polword{truć}{to poison} → \polword{zatruć}{to poison (completely)}.
  \item[\textbf{po-}] Often signals completion or doing something for a while: 
    \polword{czekać}{to wait} → \polword{poczekać}{to wait (for a bit / done waiting)}.
  \item[\textbf{na-}] Completion, sufficiency: 
    \polword{pisać}{to write} → \polword{napisać}{to write (and finish)}.
  \item[\textbf{od-}] Away from, completing a departure: 
    \polword{iść}{to go} → \polword{odejść}{to go away, depart}.
  \item[\textbf{wy-}] Out, thoroughly: 
    \polword{kopać}{to dig} → \polword{wykopać}{to dig out/up}.
  \item[\textbf{z- / s-}] Completion, bringing together: 
    \polword{bliski}{close} → \polword{zbliżyć się}{to draw closer}.
\end{description}

% ────────────────────────────────────────────────────────────
\section{Verb Conjugation}

Polish verbs encode person and number in their endings, making subject 
pronouns optional (as in Latin and Russian). The pronouns are used 
for emphasis or contrast.

\subsection*{Present tense — example: \textit{czekać} (to wait, imperfective)}

\begin{tabular}{lll}
  \textbf{Person} & \textbf{Singular} & \textbf{Plural} \\
  1st & czekam (I wait) & czekamy (we wait) \\
  2nd & czekasz (you wait) & czekacie (you wait) \\
  3rd & czeka (he/she/it waits) & czekają (they wait) \\
\end{tabular}

\subsection*{Past tense — gender marking}

Uniquely among European languages, Polish past tense verbs agree with 
the gender of the subject. This gender marking is a suffix added to the 
past stem:

\begin{description}
  \item[Masculine singular] \textbf{-ł}: \textit{powiedział} (he said), 
    \textit{spotkałem} (I [masc.] met)
  \item[Feminine singular] \textbf{-ła}: \textit{czekała} (she waited), 
    \textit{musiała} (she had to)
  \item[Neuter singular] \textbf{-ło}: \textit{miało} (it was supposed to)
  \item[Virile plural (masc. personal)] \textbf{-li}: \textit{zdecydowaliśmy} 
    (we [men/mixed] decided), \textit{ustawili się} (they positioned themselves)
  \item[Non-virile plural] \textbf{-ły}: used for groups of women or non-persons
\end{description}

\grambox{Virile vs.\ non-virile}{Polish distinguishes between groups 
containing at least one man (\textit{virile}, personal masculine plural) 
and all other plurals. This distinction affects verb endings, adjective 
forms, and pronouns throughout the grammar.}

% ────────────────────────────────────────────────────────────
\section{Key Impersonal Constructions}

Polish uses several impersonal constructions that have no direct English 
equivalent but will be encountered on virtually every page.

\begin{description}
  \item[\polword{trzeba}{one must, it is necessary}] Followed by an 
    infinitive. No subject. \textit{Trzeba kontynuować} — one must continue. 
    Equivalent to French \textit{il faut}.
  \item[\polword{można}{one can, it is possible}] Followed by an 
    infinitive. \textit{Można było tam dostać} — one could get there. 
    Equivalent to French \textit{on peut}.
  \item[\polword{warto}{it is worth}] Followed by an infinitive. 
    \textit{Warto wiedzieć} — it is worth knowing.
\end{description}

% ────────────────────────────────────────────────────────────
\section{Essential Particles and Connectives}

A small set of particles carries enormous communicative weight in Polish 
and will be encountered constantly.

\begin{description}
  \item[\polword{już}{already, now, anymore}] Marks a transition or 
    change of state. \textit{Już rozumiem} — I understand now (I've crossed 
    into understanding). With negation: \textit{już nie} — not anymore.
  \item[\polword{jeszcze}{still, yet, even more}] Marks continuation 
    or addition. \textit{Jeszcze jeden} — yet one more. With negation: 
    \textit{jeszcze nie} — not yet.
  \item[\polword{już nie}{not anymore}] Something was happening but has stopped.
  \item[\polword{jeszcze nie}{not yet}] Something has not happened but may.
  \item[\polword{też / także}{also, too}] \textit{Też} is conversational; 
    \textit{także} is slightly more formal.
  \item[\polword{wcale nie}{not at all}] Emphatic negation. 
    \textit{Wcale nie odchodzi} — it does not go away at all.
  \item[\polword{właśnie}{exactly, just now, precisely}] 
    Confirms or emphasises: \textit{właśnie} as a standalone response 
    means ``exactly'' or ``quite so.''
  \item[\polword{chyba}{probably, I suppose, surely}] Expresses 
    uncertainty or mild assertion. Very common in speech.
  \item[\polword{przecież}{after all, but surely}] Implies that 
    something should be obvious. \textit{Przecież wiesz} — you know 
    this perfectly well (after all).
\end{description}

% ────────────────────────────────────────────────────────────
\section{Conjunctions: Because}

Three words translate as ``because'' with different registers:

\begin{description}
  \item[\polword{bo}{because}] Casual, spoken, everyday. 
    The default in conversation.
  \item[\polword{ponieważ}{because}] Standard written and formal spoken. 
    Neutral and widely applicable.
  \item[\polword{gdyż}{because/for}] Literary and elevated. 
    Signals deliberate, crafted prose. Equivalent to French \textit{car}.
\end{description}

% ────────────────────────────────────────────────────────────
\section{Word Order}

Polish word order is considerably freer than English, because case 
endings — not position — identify grammatical role. However, order 
is far from arbitrary: it encodes \textbf{information structure}.

\begin{description}
  \item[New information] tends to come at the \textbf{end} of the sentence 
    (the ``rheme'' or focus position). This is why Polish often places the 
    subject after the verb, or saves the most concrete noun for last.
  \item[Given / known information] tends to come early.
  \item[Verb-first] is used for dramatic effect, as in: 
    \textit{Umarł Fiszer} — Fiszer died (the death before the name).
\end{description}

This freedom means that close translation often fails to capture 
the author's emphasis. Notes in the chapter analyses point out 
significant word-order choices where they carry meaning.

% ────────────────────────────────────────────────────────────
\section{Reflexive Verbs and \textit{się}}

The particle \textbf{się} is one of the most versatile elements in 
Polish. It functions as:

\begin{description}
  \item[Reflexive marker] The action reflects back on the subject: 
    \textit{zbliżyć się} — to bring oneself closer, to approach.
  \item[Reciprocal marker] \textit{spotkać się} — to meet (each other).
  \item[Impersonal/passive marker] \textit{mówi się} — one says, 
    it is said. Removes the agent entirely.
  \item[Part of the verb's inherent meaning] Some verbs simply 
    require \textit{się} with no separable reflexive meaning: 
    \textit{bać się} (to be afraid), \textit{cieszyć się} (to be glad).
\end{description}

% ────────────────────────────────────────────────────────────
\section{Participles}

Polish uses participial constructions extensively — more so than 
conversational English, but comparable to written French or Latin.

\begin{description}
  \item[Present active participle (\textit{imiesłów przymiotnikowy czynny})] 
    Describes an ongoing action of the noun it modifies. 
    Formed with suffix \textbf{-ący/-ąca/-ące}: 
    \textit{grupka \textbf{stojąca} przy grobie} — the group standing by the grave.
  \item[Past passive participle (\textit{imiesłów przymiotnikowy bierny})] 
    Describes a completed action done to the noun. 
    Formed with suffix \textbf{-ny/-na/-ne} or \textbf{-ty/-ta/-te}: 
    \textit{świeżo \textbf{wykopany} grób} — the freshly dug grave; 
    \textit{zapomniana epidemia} — the forgotten epidemic.
  \item[Adverbial participle (\textit{imiesłów przysłówkowy})] 
    Describes a simultaneous action of the subject. Does not decline. 
    Formed with suffix \textbf{-ąc}: 
    \textit{\textbf{odkręcając} nakrętkę} — while unscrewing the cap.
\end{description}
